\chapter{Minimum Spanning Trees}

In Part~I, we traversed unweighted graphs where every edge counted as a single
step. In Part~II, edges acquire real-valued weights -- representing costs,
distances, or capacities. When weights measure the cost of constructing or
maintaining connections, our goal is to span all $n$ vertices using as few edges
as possible while minimizing the total cost. Any minimal connected subgraph is a
spanning tree with exactly $n-1$ edges. This chapter develops the minimum spanning
tree (MST) problem from the cut and cycle properties, showing why
greedy choices by Kruskal, Prim, and Bor\r{u}vka are provably optimal.

\chapterline

\section{The problem and a picture}

\begin{definition}[Spanning tree, MST]
Let $\Gph=(V,E)$ be a connected undirected graph with a weight $w(e)\in\R$ on each edge. A \emph{spanning tree} is an acyclic connected subgraph $T=(V,E_T)$ with $E_T\subseteq E$; necessarily $\abs{E_T}=n-1$ where $n=\abs V$. Its weight is $w(T)=\sum_{e\in E_T} w(e)$. A \emph{minimum spanning tree} is a spanning tree of least total weight. If $\Gph$ is disconnected, the analogous object is a \emph{minimum spanning forest}, one tree per connected component.
\end{definition}

\begin{intuitionbox}
Think of the vertices as towns and the edges as candidate stretches of fibre, each labelled by what it costs to dig. You must be able to route a signal between any two towns, and you are paying by the metre. You do not care about path lengths or redundancy -- only that the network is connected and that the total bill is as small as possible. The answer is always a tree: any cycle contains an edge you could delete while staying connected, and deleting a positive-cost edge never hurts. The question is \emph{which} tree.
\end{intuitionbox}

Here is the running example we will trace by hand throughout the chapter.

\begin{figure}[H]
\centering
\begin{tikzpicture}[
  v/.style={circle,draw=nodeblue,thick,fill=blue!4,minimum size=8mm,inner sep=0pt,font=\small},
  edgelabel/.style={font=\footnotesize,inner sep=1.5pt},
  scale=1.0]
\node[v] (A) at (0,2) {$A$};
\node[v] (B) at (3,3) {$B$};
\node[v] (C) at (6,2.4) {$C$};
\node[v] (D) at (1.2,0) {$D$};
\node[v] (E) at (4,0.2) {$E$};
\node[v] (F) at (6.6,0.4) {$F$};
\draw (A) -- node[edgelabel,above left]{$4$} (B);
\draw (A) -- node[edgelabel,below left]{$1$} (D);
\draw (B) -- node[edgelabel,above left]{$2$} (D);
\draw (B) -- node[edgelabel,above]{$6$} (C);
\draw (B) -- node[edgelabel,right=1pt]{$5$} (E);
\draw (D) -- node[edgelabel,below]{$3$} (E);
\draw (C) -- node[edgelabel,above left]{$7$} (E);
\draw (C) -- node[edgelabel,right=1pt]{$8$} (F);
\draw (E) -- node[edgelabel,below]{$9$} (F);
\end{tikzpicture}
\caption{The running graph $\Gph_0$. Six vertices, nine weighted edges, all weights distinct.}
\label{ch06:running-graph}
\end{figure}

We will keep all weights distinct in $\Gph_0$ on purpose: distinct weights force a \emph{unique} MST, which makes the traces unambiguous. Section~\ref{ch06:uniqueness} handles ties. For the record, the MST of $\Gph_0$ turns out to be $\set{AD(1),\,BD(2),\,DE(3),\,BC(6),\,CF(8)}$ with total weight $20$; every algorithm below will rediscover it.

\section{Two structural laws: the cut and the cycle}

Every correct MST algorithm rests on two facts about which edges \emph{must} or \emph{must not} appear. They are duals of each other: one is phrased over cuts, the other over cycles.

\begin{definition}[Cut]
A \emph{cut} $(S,\overline S)$ is a partition of $V$ into a nonempty set $S$ and its complement $\overline S=V\setminus S$, both nonempty. An edge \emph{crosses} the cut if it has one endpoint in $S$ and one in $\overline S$. A cut \emph{respects} a set of edges $F$ if no edge of $F$ crosses it.
\end{definition}

\begin{focusbox}
\textbf{Cut property.} For any cut, the cheapest crossing edge can be safely put into an MST. \\
\textbf{Cycle property.} For any cycle, the most expensive edge on it can be safely left out of every MST. \\
Both ``safely'' statements become \emph{exactly one MST forced} when the relevant minimum/maximum is unique.
\end{focusbox}

\begin{theorem}[Cut property]
\label{ch06:cut-property}
Let $F\subseteq E$ be a subset of edges contained in some MST of $\Gph$. Let $(S,\overline S)$ be any cut that respects $F$, and let $e$ be a crossing edge of minimum weight among all crossing edges. Then $F\cup\set e$ is also contained in some MST.
\end{theorem}

\begin{whybox}
The proof is the \emph{exchange argument} (Figure~\ref{ch06:cut-exchange}), the standard way to prove that a greedy choice is optimal. Take any MST $T$ that contains $F$. If it already uses $e$, we are done. If not, adding $e$ to $T$ creates exactly one cycle; that cycle must cross the cut a second time (it left $S$ via $e$, it has to come back), so it contains another crossing edge $e'$. Swap them: remove $e'$, add $e$. The result is still a spanning tree, it still contains $F$ (because the cut respected $F$, so $e'\notin F$), and since $e$ was a cheapest crossing edge, $w(e)\le w(e')$, so the swap did not increase the weight. A tree no heavier than an MST is an MST.
\end{whybox}

\begin{proof}
Let $T$ be an MST with $F\subseteq E_T$. If $e\in E_T$ then $F\cup\set e\subseteq E_T$ and we are finished. Otherwise $T+e$ contains a unique cycle $\gamma$ (adding one edge to a tree always creates exactly one cycle). The endpoints of $e$ lie on opposite sides of the cut $(S,\overline S)$, so traversing $\gamma$ we cross from $S$ to $\overline S$ at $e$ and must cross back at least once; pick such a return edge $e'\in\gamma$, $e'\ne e$, with $e'$ crossing the cut. Because the cut respects $F$ and $e'$ crosses it, $e'\notin F$.

Form $T'=T-e'+e$. Removing $e'$ from $T+e$ breaks the only cycle, so $T'$ is acyclic; it has $n-1$ edges and is acyclic, hence a spanning tree. Its weight is $w(T')=w(T)-w(e')+w(e)$. Since $e$ is a minimum-weight crossing edge and $e'$ also crosses, $w(e)\le w(e')$, so $w(T')\le w(T)$. But $T$ is minimum, so $w(T')=w(T)$ and $T'$ is also an MST. Finally $F\cup\set e\subseteq E_{T'}$: every edge of $F$ survived (we only deleted $e'\notin F$) and $e$ was added.
\end{proof}

\begin{corollary}[Cut property, sharp form]
\label{ch06:cut-sharp}
If, in addition, $e$ is the \emph{unique} minimum-weight crossing edge of the cut, then \emph{every} MST containing $F$ contains $e$.
\end{corollary}

\begin{proof}
Suppose some MST $T\supseteq F$ omits $e$. Run the exchange as above to get $T'=T-e'+e$ with $w(e)<w(e')$ (strict, by uniqueness), so $w(T')<w(T)$, contradicting minimality of $T$.
\end{proof}

\begin{figure}[H]
\centering
\begin{tikzpicture}[
  v/.style={circle,draw=nodeblue,thick,fill=blue!4,minimum size=7mm,inner sep=0pt,font=\small},
  edgelabel/.style={font=\footnotesize,fill=white,inner sep=1.5pt},
  treeedge/.style={draw=edgegray,line width=1pt},
  newedge/.style={draw=targetgreen,line width=1.6pt},
  swapedge/.style={draw=trapred,line width=1.6pt,dashed},
  scale=1.0]
% left side S
\node[v] (s1) at (0,1.6) {};
\node[v] (s2) at (0.2,0) {};
\node[v] (u)  at (1.8,0.8) {$u$};
\draw[treeedge] (s1) -- (u);
\draw[treeedge] (s2) -- (u);
% right side Sbar
\node[v] (v)  at (4.4,1.6) {$v$};
\node[v] (t1) at (6.2,1.6) {};
\node[v] (t2) at (5.4,0) {};
\draw[treeedge] (v) -- (t1);
\draw[treeedge] (v) -- (t2);
% the cut line (in background so edge labels stay crisp)
\draw[edgegray,dashed,line width=0.8pt] (3.1,2.4) -- (3.1,-0.6);
% the cheap crossing edge e = (u,v)
\draw[newedge] (u) -- node[pos=0.35, above=2pt, edgelabel, text=targetgreen!85!black]{$e$} (v);
% the second crossing edge e' on the cycle, here (s1)..(t1) type
\draw[swapedge] (s1) to[bend left=20] node[pos=0.35, above=2pt, edgelabel, text=trapred!85!black]{$e'$} (t1);
\node[font=\footnotesize] at (0.9,2.3) {$S$};
\node[font=\footnotesize] at (5.3,2.3) {$\overline S$};
\end{tikzpicture}
\caption{The exchange argument for the cut property. The cut (dashed vertical line) splits the tree; $e$ (green) is the cheapest crossing edge. Adding $e$ to the tree closes a cycle that must cross the cut a second time at some tree edge $e'$ (red dashed). Since $w(e)\le w(e')$, swapping $e'$ out for $e$ does not increase the weight.}
\label{ch06:cut-exchange}
\end{figure}

The cycle property is the mirror image and is proved the same way, exchanging in the other direction.

\begin{theorem}[Cycle property]
\label{ch06:cycle-property}
Let $\gamma$ be any cycle in $\Gph$ and let $f$ be an edge of strictly maximum weight on $\gamma$. Then $f$ belongs to \emph{no} MST.
\end{theorem}

\begin{proof}
Suppose an MST $T$ contains $f$. Deleting $f$ splits $T$ into two components, inducing a cut $(S,\overline S)$ that $f$ crosses. The cycle $\gamma$ enters and leaves $S$, so it contains another crossing edge $g\ne f$ with $g\notin E_T$ (the only $T$-edge crossing this cut is $f$). By the choice of $f$, $w(g)<w(f)$. Then $T-f+g$ is a spanning tree of weight $w(T)-w(f)+w(g)<w(T)$, contradicting minimality.
\end{proof}

The two panels below summarize this duality between including the lightest edge across a cut and excluding the heaviest edge along a cycle.

\begin{center}
\begin{tikzpicture}[
    >=Stealth,
    vnode/.style={circle, draw=sourceblue, fill=sourceblue!10, thick, minimum size=7.5mm, font=\small\bfseries},
    cutedge/.style={draw=edgegray, thick},
    lightedge/.style={draw=targetgreen!85!black, very thick},
    heavyedge/.style={draw=trapred!90!black, very thick, dashed},
    paneltitle/.style={font=\footnotesize\scshape, text=gray!85!black}
]
  % Panel 1: Cut Property (Light edge safe to add)
  \begin{scope}[xshift=0cm]
    \node[paneltitle] at (2.2, 3.2) {Cut Property: Lightest crossing edge included};

    % Cut boundary regions
    \draw[draw=sourceblue!50, fill=sourceblue!5, rounded corners=6pt, dashed] (-0.6, -0.6) rectangle (1.4, 2.6);
    \draw[draw=targetgreen!50, fill=targetgreen!5, rounded corners=6pt, dashed] (3.0, -0.6) rectangle (5.0, 2.6);
    \node[font=\footnotesize\bfseries, text=sourceblue!80!black] at (0.4, 2.3) {Set $S$};
    \node[font=\footnotesize\bfseries, text=targetgreen!80!black] at (4.0, 2.3) {Set $\overline{S}$};

    \node[vnode] (s1) at (0.4, 1.5) {$u_1$};
    \node[vnode] (s2) at (0.4, 0.2) {$u_2$};
    \node[vnode] (t1) at (4.0, 1.5) {$v_1$};
    \node[vnode] (t2) at (4.0, 0.2) {$v_2$};

    \draw[cutedge] (s1) -- node[midway, above, font=\scriptsize] {$8$} (t1);
    \draw[lightedge] (s1) -- node[midway, sloped, above=1pt, font=\scriptsize\bfseries, text=targetgreen!90!black] {$\mathbf{4}$ (lightest)} (t2);
    \draw[cutedge] (s2) -- node[midway, below, font=\scriptsize] {$6$} (t2);

    \node[anchor=center, font=\scriptsize, text=targetgreen!90!black] at (2.2, -1.0)
      {Cheapest crossing edge $w=4$ safe to add};
  \end{scope}

  % Panel 2: Cycle Property (Heaviest edge safe to delete)
  \begin{scope}[xshift=7.2cm]
    \node[paneltitle] at (2.0, 3.2) {Cycle Property: Heaviest edge excluded};

    \node[vnode] (c1) at (0.6, 1.8) {$a$};
    \node[vnode] (c2) at (3.4, 1.8) {$b$};
    \node[vnode] (c3) at (3.4, 0.0) {$c$};
    \node[vnode] (c4) at (0.6, 0.0) {$d$};

    \draw[cutedge] (c1) -- node[midway, above, font=\scriptsize] {$2$} (c2);
    \draw[cutedge] (c2) -- node[midway, right, font=\scriptsize] {$5$} (c3);
    \draw[cutedge] (c3) -- node[midway, below, font=\scriptsize] {$3$} (c4);
    \draw[heavyedge] (c1) -- node[midway, left, font=\scriptsize\bfseries, text=trapred!90!black] {$\mathbf{9}$ (heaviest)} (c4);

    \node[anchor=center, font=\scriptsize, text=trapred!90!black] at (2.0, -1.0)
      {Heaviest cycle edge $w=9$ excluded from MST};
  \end{scope}
\end{tikzpicture}
\end{center}

\begin{intuitionbox}
The cut property tells you what to \emph{include}; the cycle property tells you what to \emph{exclude}. Kruskal and Boruvka are most naturally read as repeated applications of the cut property building the tree up; the cycle property is the tool of choice when you want to argue an edge is \emph{wrong} -- which is the basis of MST verification in Section~\ref{ch06:verify}. The reverse-delete algorithm (sort descending, delete an edge whenever doing so keeps the graph connected) turns the cycle property directly into an algorithm.
\end{intuitionbox}

\begin{trapbox}
The cut property does not say ``a minimum crossing edge is in \emph{every} MST.'' Only the \emph{sharp} form, with a \emph{unique} minimum, forces the edge. With ties, a cheapest crossing edge is in \emph{some} MST, which is exactly what a greedy algorithm needs -- but two different tie-breaks can build two different trees. Conflating these two readings is the single most common MST proof bug.
\end{trapbox}

\section{The generic greedy algorithm}

All three classical algorithms are instances of one meta-algorithm: maintain a forest $F$ that is a subset of \emph{some} MST (we call such an $F$ \emph{promising}), and repeatedly add a \emph{safe} edge -- a minimum-weight edge crossing some cut that respects $F$ -- until $F$ spans.

\begin{algorithm}[H]
\caption{Generic-MST -- the greedy meta-algorithm}
\label{ch06:generic}
\KwIn{connected weighted graph $\Gph=(V,E,w)$}
\KwOut{edge set of an MST}
$F\gets\emptyset$\;
\While{$F$ does not span $\Gph$}{
  pick a cut $(S,\overline S)$ that respects $F$\;
  let $e$ be a minimum-weight edge crossing $(S,\overline S)$\;
  $F\gets F\cup\set e$\;
}
\Return $F$\;
\end{algorithm}

\begin{whybox}
\textbf{Correctness of the meta-algorithm.} The invariant ``$F$ is promising'' holds at the start ($\emptyset$ is contained in every MST) and is preserved by every iteration: the chosen cut respects $F$ and $e$ is a minimum crossing edge, so Theorem~\ref{ch06:cut-property} keeps $F\cup\set e$ promising. The loop adds one edge per iteration without ever creating a cycle (a safe edge crosses a respecting cut, so it joins two different trees of $F$), so after $n-1$ iterations $F$ is a spanning tree that is promising -- i.e.\ contained in an MST -- and a spanning tree contained in an MST \emph{is} that MST. Every concrete algorithm below differs only in \emph{which cut it picks} and \emph{how it finds the cheapest crossing edge fast}.
\end{whybox}

The three classical specializations:

\begin{itemize}
\item \textbf{Kruskal}: consider edges globally cheapest-first; each accepted edge is safe for the cut separating the two components it joins.
\item \textbf{Prim}: keep one growing tree $S$; repeatedly add the cheapest edge leaving $S$.
\item \textbf{Boruvka}: in parallel, every current component picks its own cheapest outgoing edge; add them all at once.
\end{itemize}

\section{Kruskal's algorithm}

\begin{historybox}
The minimum spanning tree problem has a long history. Otakar Bor\r{u}vka formulated the problem and discovered the first algorithm in 1926 to design an economical electrical grid for rural Moravia. Vojt\v{e}ch Jarn\'{\i}k discovered the growing-tree method in 1930, later rediscovered independently by Robert Prim in 1957 and Edsger Dijkstra in 1959. Joseph Kruskal proposed the global sorting method in 1956, which became near-linear when paired with the disjoint-set forest (Galler and Fischer, 1964) analyzed by Tarjan in 1975. Later work improved the bounds: Karger, Klein, and Tarjan (1995) achieved randomized linear time $\Oh(m)$, while whether a deterministic comparison-based algorithm can solve the problem in strictly linear time remains one of the deepest open questions in graph algorithms.
\end{historybox}

\begin{intuitionbox}
Sort the edges from cheapest to dearest and walk down the list. Take each edge unless it would close a cycle, i.e.\ unless its two endpoints are already connected by edges you took earlier. You are greedily buying the cheapest edge that buys you something -- that merges two as-yet-separate pieces of the network.
\end{intuitionbox}

The cycle test ``are these two endpoints already connected?'' is exactly a \emph{disjoint-set} (union--find) query: maintain the connected components of the chosen-so-far forest, with \texttt{find}$(v)$ returning a canonical representative and \texttt{union} merging two components.

\begin{algorithm}[H]
\caption{Kruskal-MST}
\label{ch06:kruskal}
\KwIn{connected weighted graph $\Gph=(V,E,w)$}
\KwOut{MST edge set $T$}
$T\gets\emptyset$\;
\lForEach{$v\in V$}{\textsc{Make-Set}$(v)$}
sort $E$ so that $w(e_1)\le w(e_2)\le\dots\le w(e_m)$\;
\ForEach{edge $e=(u,v)$ in nondecreasing weight order}{
  \If{\textsc{Find}$(u)\ne$ \textsc{Find}$(v)$}{
    $T\gets T\cup\set e$\;
    \textsc{Union}$(u,v)$\;
    \lIf{$\abs T=n-1$}{\textbf{break}}
  }
}
\Return $T$\;
\end{algorithm}

\subsection{Union--find in one paragraph}

Represent each set as a rooted tree of parent pointers; the root is the representative. \texttt{find} walks to the root, and during \emph{path compression} repoints every node it passes directly at the root. \texttt{union} links the two roots, attaching the shorter tree under the taller (\emph{union by rank}). With both heuristics, any sequence of $m$ operations on $n$ elements costs $\Oh(m\,\alpha(n))$, where $\alpha$ is the inverse Ackermann function -- below $5$ for every $n$ writable in this universe. We treat union--find as a black box costing $\Oh(\alpha(n))$ amortized per operation here; the data structure and its analysis get their own chapter.

\begin{labbox}
Trace Kruskal on $\Gph_0$ (Figure~\ref{ch06:running-graph}). Sorted edges:
$AD(1)$, $BD(2)$, $DE(3)$, $AB(4)$, $BE(5)$, $BC(6)$, $CE(7)$, $CF(8)$, $EF(9)$.

\begin{center}
\begin{tabular}{clll}
\toprule
edge & weight & action & components after \\
\midrule
$AD$ & $1$ & accept & $\set{A,D}\,\set B\,\set C\,\set E\,\set F$ \\
$BD$ & $2$ & accept & $\set{A,B,D}\,\set C\,\set E\,\set F$ \\
$DE$ & $3$ & accept & $\set{A,B,D,E}\,\set C\,\set F$ \\
$AB$ & $4$ & \emph{reject} ($A,B$ joined) & unchanged \\
$BE$ & $5$ & \emph{reject} ($B,E$ joined) & unchanged \\
$BC$ & $6$ & accept & $\set{A,B,C,D,E}\,\set F$ \\
$CE$ & $7$ & \emph{reject} ($C,E$ joined) & unchanged \\
$CF$ & $8$ & accept & $\set{A,B,C,D,E,F}$ -- spans, stop \\
\bottomrule
\end{tabular}
\end{center}

Result $\set{AD,BD,DE,BC,CF}$, weight $1+2+3+6+8=20$. Note how every \emph{rejected} edge is the heaviest on the cycle it would have closed: $AB(4)$ on $A\!-\!D\!-\!B\!-\!A$, $CE(7)$ on $C\!-\!B\!-\!D\!-\!E\!-\!C$. That is the cycle property in action.
\end{labbox}

\begin{figure}[H]
\centering
\begin{tikzpicture}[
  v/.style={circle,draw=nodeblue,thick,fill=blue!4,minimum size=8mm,inner sep=0pt,font=\small},
  edgelabel/.style={font=\footnotesize,inner sep=1.5pt},
  tree/.style={draw=targetgreen,line width=1.4pt},
  cut/.style={draw=edgegray,dashed,line width=0.8pt},
  scale=1.0]
\node[v] (A) at (0,2) {$A$};
\node[v] (B) at (3,3) {$B$};
\node[v] (C) at (6,2.4) {$C$};
\node[v] (D) at (1.2,0) {$D$};
\node[v] (E) at (4,0.2) {$E$};
\node[v] (F) at (6.6,0.4) {$F$};
\draw[cut] (A) -- node[edgelabel,above left]{$4$} (B);
\draw[tree] (A) -- node[edgelabel,below left]{$1$} (D);
\draw[tree] (B) -- node[edgelabel,above left]{$2$} (D);
\draw[tree] (B) -- node[edgelabel,above]{$6$} (C);
\draw[cut] (B) -- node[edgelabel,right=1pt]{$5$} (E);
\draw[tree] (D) -- node[edgelabel,below]{$3$} (E);
\draw[cut] (C) -- node[edgelabel,above left]{$7$} (E);
\draw[tree] (C) -- node[edgelabel,right=1pt]{$8$} (F);
\draw[cut] (E) -- node[edgelabel,below]{$9$} (F);
\end{tikzpicture}
\caption{The MST of $\Gph_0$ in green ($AD,BD,DE,BC,CF$), total weight $20$. Dashed edges are excluded; each is the maximum on some cycle.}
\label{ch06:mst}
\end{figure}

\begin{whybox}
\textbf{Why Kruskal is correct.} Each accepted edge $e=(u,v)$ joins two distinct trees of the current forest $F$. Take $S$ to be the vertex set of $u$'s tree. The cut $(S,\overline S)$ respects $F$ (no forest edge crosses it -- $S$ is a whole component). Every still-unprocessed edge is at least as heavy as $e$ (we go in sorted order), and every \emph{already}-processed crossing edge was rejected, which only happens inside a single component, so no processed edge crosses $(S,\overline S)$ at all. Hence $e$ is a minimum-weight crossing edge, and the cut property makes it safe. Kruskal is Generic-MST with the cut chosen retroactively around each accepted edge.
\end{whybox}

\begin{complexitybox}
Sorting $m$ edges: $\Oh(m\log m)=\Oh(m\log n)$ (since $m\le n^2$, $\log m\le 2\log n$). The union--find work: $2m$ finds and $n-1$ unions, $\Oh(m\,\alpha(n))$. Total $\Oh(m\log n)$, dominated by the sort. If the weights are small integers and can be radix-sorted in $\Oh(m)$, Kruskal runs in $\Oh(m\,\alpha(n))$ -- essentially linear.
\end{complexitybox}

\section{Prim's algorithm}

Where Kruskal's algorithm grows a forest across the entire graph and merges components in weight order, Prim's algorithm (Jarn\'{\i}k, 1930; Prim, 1957; Dijkstra, 1959) grows one connected tree outward from a single starting vertex. Its loop mirrors Dijkstra's shortest-path algorithm, except each outside vertex is prioritized by the weight of its single lightest edge into the tree rather than its accumulated path distance from the root.

\begin{intuitionbox}
Grow a single tree like a crystal. Start from any vertex. At each step look at all edges with exactly one endpoint already in the tree -- the \emph{frontier} -- and pull in the cheapest one, absorbing its outside endpoint. Repeat until everyone is inside. Prim always picks the cut $(S,\overline S)$ where $S$ is the current tree, so safety is immediate from the cut property.
\end{intuitionbox}

The efficiency question is: how do we repeatedly find the cheapest frontier edge? Keep, for every vertex $v\notin S$, the cheapest known edge connecting it to $S$, call it $\mathrm{key}[v]$, in a priority queue. When a vertex is absorbed, relax its neighbours' keys, exactly as Dijkstra relaxes distances.

\begin{algorithm}[H]
\caption{Prim-MST (binary-heap version)}
\label{ch06:prim}
\KwIn{connected weighted graph $\Gph=(V,E,w)$, start vertex $r$}
\KwOut{MST given by parent pointers $\pi$}
\lForEach{$v\in V$}{$\mathrm{key}[v]\gets+\infty$;\ $\pi[v]\gets\textsc{nil}$;\ $v$ \textbf{unmarked}}
$\mathrm{key}[r]\gets 0$\;
priority queue $Q\gets V$ keyed by $\mathrm{key}$\;
\While{$Q\ne\emptyset$}{
  $u\gets\textsc{Extract-Min}(Q)$;\ \ mark $u$\;
  \ForEach{edge $(u,v)$ with $v\in Q$ unmarked}{
    \If{$w(u,v)<\mathrm{key}[v]$}{
      $\mathrm{key}[v]\gets w(u,v)$;\ \ $\pi[v]\gets u$;\ \ \textsc{Decrease-Key}$(Q,v)$\;
    }
  }
}
\Return $\set{(v,\pi[v]):v\ne r}$\;
\end{algorithm}

\begin{labbox}
Trace Prim on $\Gph_0$ from $r=A$. Columns are $\mathrm{key}[\cdot]$; $\bullet$ marks the extracted (absorbed) vertex; $-$ means already in the tree; the boldface entry in each row is the minimum key, i.e.\ the vertex absorbed in the next row.

\begin{center}
\begin{tabular}{lcccccc}
\toprule
step (absorb) & $A$ & $B$ & $C$ & $D$ & $E$ & $F$ \\
\midrule
init                & $0$ & $\infty$ & $\infty$ & $\infty$ & $\infty$ & $\infty$ \\
$A^\bullet$         & $-$ & $4$ & $\infty$ & $\mathbf 1$ & $\infty$ & $\infty$ \\
$D^\bullet$ (via $A$)& $-$ & $\mathbf 2$ & $\infty$ & $-$ & $3$ & $\infty$ \\
$B^\bullet$ (via $D$)& $-$ & $-$ & $6$ & $-$ & $\mathbf 3$ & $\infty$ \\
$E^\bullet$ (via $D$)& $-$ & $-$ & $\mathbf 6$ & $-$ & $-$ & $9$ \\
$C^\bullet$ (via $B$)& $-$ & $-$ & $-$ & $-$ & $-$ & $\mathbf 8$ \\
$F^\bullet$ (via $C$)& $-$ & $-$ & $-$ & $-$ & $-$ & $-$ \\
\bottomrule
\end{tabular}
\end{center}

The parent edges chosen, in order: $AD(1),DB(2),DE(3),BC(6),CF(8)$ -- the same tree, weight $20$. Watch the key of $C$: it first becomes $6$ when $B$ is absorbed, and $E$ later offers $C$ a connection of weight $7$, which is rejected because $7>6$. Watch $F$: it is offered $9$ via $E$, then improved to $8$ via $C$.
\end{labbox}

The snapshot below captures Prim's state right after absorbing $S=\{A,D,B\}$: the cut $(S,\overline S)$ separates the growing tree from the outside vertices, whose priority-queue keys record each vertex's lightest edge into $S$.

\begin{center}
\begin{tikzpicture}[
  vtree/.style={circle,draw=sourceblue!90!black,fill=sourceblue!22,very thick,minimum size=8mm,inner sep=0pt,font=\small\bfseries},
  vfront/.style={circle,draw=quotientorange!90!black,fill=quotientorange!15,very thick,minimum size=8mm,inner sep=0pt,font=\small\bfseries},
  vfar/.style={circle,draw=gray!60,fill=gray!8,thick,minimum size=8mm,inner sep=0pt,font=\small},
  edgelabel/.style={font=\footnotesize,inner sep=1.5pt},
  tree/.style={draw=sourceblue!85!black,line width=1.6pt},
  bestcut/.style={draw=targetgreen!85!black,line width=1.8pt},
  candcut/.style={draw=quotientorange!85!black,line width=1.2pt,dashed},
  stale/.style={draw=edgegray,line width=0.8pt,dotted},
  scale=1.0]
  \begin{scope}[on background layer]
    \draw[fill=sourceblue!6,draw=sourceblue!45,dashed,rounded corners=10pt]
      (-0.65,-0.65) -- (3.65,-0.65) -- (3.65,3.85) -- (-0.65,3.85) -- cycle;
    \node[font=\scriptsize\bfseries,text=sourceblue!85!black] at (1.1,3.55) {Absorbed Tree $S=\{A,D,B\}$};
  \end{scope}
  \node[vtree] (A) at (0,2) {$A$};
  \node[vtree] (B) at (3,3) {$B$};
  \node[vfront] (C) at (6.2,2.4) {$C$};
  \node[vtree] (D) at (1.2,0) {$D$};
  \node[vfront] (E) at (4.6,0.2) {$E$};
  \node[vfar] (F) at (7.4,0.4) {$F$};

  \node[font=\scriptsize\bfseries,text=targetgreen!80!black,below=2pt of E] {$\mathrm{key}[E]=\mathbf{3}$ (next)};
  \node[font=\scriptsize\bfseries,text=quotientorange!90!black,above=2pt of C] {$\mathrm{key}[C]=6$};
  \node[font=\scriptsize,text=gray!70!black,below=2pt of F] {$\mathrm{key}[F]=\infty$};

  \draw[stale] (A) -- node[edgelabel,above left]{$4$} (B);
  \draw[tree] (A) -- node[edgelabel,below left]{$1$} (D);
  \draw[tree] (B) -- node[edgelabel,above left]{$2$} (D);
  \draw[candcut] (B) -- node[edgelabel,above]{$6$} (C);
  \draw[stale] (B) -- node[edgelabel,right=1pt]{$5$} (E);
  \draw[bestcut] (D) -- node[edgelabel,below,text=targetgreen!85!black]{$\mathbf{3}$} (E);
  \draw[stale] (C) -- node[edgelabel,above left]{$7$} (E);
  \draw[stale] (C) -- node[edgelabel,right=1pt]{$8$} (F);
  \draw[stale] (E) -- node[edgelabel,below]{$9$} (F);
\end{tikzpicture}
\end{center}

\begin{whybox}
\textbf{Why Prim is correct.} Invariant: the marked set $S$ together with the parent edges among them forms a tree that is promising, and for each unmarked $v$, $\mathrm{key}[v]$ is the weight of the lightest edge from $v$ into $S$ (or $\infty$). When \textsc{Extract-Min} returns $u$, $\mathrm{key}[u]$ is the minimum over \emph{all} unmarked vertices, hence the edge $(\pi[u],u)$ is a minimum-weight edge crossing the cut $(S,\overline S)$. The cut respects the current tree (S is exactly the marked component). By Theorem~\ref{ch06:cut-property}, adding it keeps the tree promising. After $n$ extractions $S=V$ and the promising tree spans, so it is an MST.
\end{whybox}

\begin{complexitybox}
With a binary heap: $n$ \textsc{Extract-Min} at $\Oh(\log n)$ each and up to $m$ \textsc{Decrease-Key} at $\Oh(\log n)$ each, total $\Oh((n+m)\log n)=\Oh(m\log n)$ for a connected graph. With a Fibonacci heap, \textsc{Decrease-Key} is $\Oh(1)$ amortized and \textsc{Extract-Min} is $\Oh(\log n)$, giving $\Oh(m+n\log n)$ -- asymptotically better for dense graphs. With a simple array (no heap), each \textsc{Extract-Min} is a $\Oh(n)$ scan and \textsc{Decrease-Key} is $\Oh(1)$, giving $\Oh(n^2)$, which is the right choice when $m=\Th(n^2)$.
\end{complexitybox}

\begin{trapbox}
Prim's priority is the \emph{edge weight} $w(u,v)$, not a running sum $\mathrm{key}[u]+w(u,v)$. If you copy Dijkstra and accumulate, you compute a shortest-path tree, not an MST -- a different and usually heavier tree. The two algorithms look almost identical in code but compute different trees; confusing them is a classic interview trap.
\end{trapbox}

\section{Boruvka's algorithm}

Bor\r{u}vka's 1926 algorithm (rediscovered by Sollin in the 1960s) grows components in parallel rounds and serves as the vertex-shrinking step in modern parallel and randomized linear-time MST solvers.

\begin{intuitionbox}
In each round, every current component looks at all edges leaving it and selects its single cheapest one. Add all of those selected edges simultaneously. Each round at least halves the number of components, because every component merges with at least one other. After $\Oh(\log n)$ rounds, one component remains: the MST.
\end{intuitionbox}

\begin{algorithm}[H]
\caption{Boruvka-MST}
\label{ch06:boruvka}
\KwIn{connected weighted graph $\Gph=(V,E,w)$ with distinct edge weights}
\KwOut{MST edge set $T$}
$T\gets\emptyset$; initialize union--find with each $v$ its own component\;
\While{number of components $>1$}{
  \ForEach{component $C$}{$\mathrm{cheap}[C]\gets\textsc{nil}$}
  \ForEach{edge $(u,v)\in E$ with \textsc{Find}$(u)\ne$\textsc{Find}$(v)$}{
    \lForEach{$C\in\set{\textsc{Find}(u),\textsc{Find}(v)}$}{
      \lIf{$\mathrm{cheap}[C]=\textsc{nil}$ \textbf{or} $w(u,v)<w(\mathrm{cheap}[C])$}{$\mathrm{cheap}[C]\gets(u,v)$}
    }
  }
  \ForEach{component $C$ with $\mathrm{cheap}[C]\ne\textsc{nil}$}{
    let $e=\mathrm{cheap}[C]=(u,v)$\;
    \If{\textsc{Find}$(u)\ne$\textsc{Find}$(v)$}{
      $T\gets T\cup\set e$;\ \ \textsc{Union}$(u,v)$\;
    }
  }
}
\Return $T$\;
\end{algorithm}

\begin{whybox}
\textbf{Why Boruvka is correct.} For a component $C$, let $e$ be its cheapest outgoing edge. Take the cut $(C,\overline C)$. It respects the forest built so far (the forest edges live inside components). $e$ is by definition a minimum-weight crossing edge of this cut, so it is safe by Theorem~\ref{ch06:cut-property}. All components apply this simultaneously; with \emph{distinct} weights no harmful conflict arises -- if $C$ and $C'$ both pick edge $e$ between them, it is added once (the \texttt{find} re-check in the merge phase prevents a duplicate or a cycle). Each chosen edge is independently safe, and the union of safe edges is acyclic because each one strictly merges two components.
\end{whybox}

\begin{trapbox}
\textbf{Why ``distinct weights'' matters here.} Suppose a $4$-cycle has all edges of weight $1$. Each component's cheapest outgoing edge is some weight-$1$ edge, and if every component independently picks the edge to its clockwise neighbour, you add all four and create a cycle. The standard fix is a deterministic tie-break: compare $(w(e),\text{edge-id})$ lexicographically, which makes weights effectively distinct and kills the symmetry. With that tie-break Boruvka is correct on arbitrary inputs.
\end{trapbox}

\begin{labbox}
Trace Boruvka on $\Gph_0$.

\emph{Round 1.} Each singleton picks its cheapest incident edge:
$A\!\to\!AD(1)$, $B\!\to\!BD(2)$, $C\!\to\!BC(6)$, $D\!\to\!AD(1)$, $E\!\to\!DE(3)$, $F\!\to\!CF(8)$.
The distinct chosen edges are $\set{AD,BD,BC,DE,CF}$ (Figure~\ref{ch06:boruvka-round}). Merging: $A,D$ join (via $AD$); $B$ joins them (via $BD$); $E$ joins (via $DE$); $C$ joins (via $BC$); $F$ joins (via $CF$). That already yields one component of all six vertices.

So $\Gph_0$ finishes in a single Boruvka round, producing $\set{AD,BD,BC,DE,CF}$, weight $20$ -- the same MST. (On graphs that do not collapse so fast, you would shrink each component to a supernode and repeat.)
\end{labbox}

\begin{figure}[H]
\centering
\begin{tikzpicture}[
  v/.style={circle,draw=nodeblue,thick,fill=blue!4,minimum size=8mm,inner sep=0pt,font=\small},
  edgelabel/.style={font=\footnotesize,inner sep=1.5pt},
  picked/.style={draw=targetgreen,line width=1.6pt},
  other/.style={draw=edgegray,line width=0.7pt},
  scale=1.0]
\node[v] (A) at (0,2) {$A$};
\node[v] (B) at (3,3) {$B$};
\node[v] (C) at (6,2.4) {$C$};
\node[v] (D) at (1.2,0) {$D$};
\node[v] (E) at (4,0.2) {$E$};
\node[v] (F) at (6.6,0.4) {$F$};
\draw[other]  (A) -- node[edgelabel,above left]{$4$} (B);
\draw[picked] (A) -- node[edgelabel,below left]{$1$} (D);
\draw[picked] (B) -- node[edgelabel,above left]{$2$} (D);
\draw[picked] (B) -- node[edgelabel,above]{$6$} (C);
\draw[other]  (B) -- node[edgelabel,right=1pt]{$5$} (E);
\draw[picked] (D) -- node[edgelabel,below]{$3$} (E);
\draw[other]  (C) -- node[edgelabel,above left]{$7$} (E);
\draw[picked] (C) -- node[edgelabel,right=1pt]{$8$} (F);
\draw[other]  (E) -- node[edgelabel,below]{$9$} (F);
\end{tikzpicture}
\caption{Boruvka, round 1 on $\Gph_0$. Each vertex's cheapest incident edge is highlighted: $A$ and $D$ both choose $AD(1)$, $B$ chooses $BD(2)$, $E$ chooses $DE(3)$, $C$ chooses $BC(6)$, $F$ chooses $CF(8)$. The five distinct chosen edges already merge all six vertices into one component, so the round finishes the MST.}
\label{ch06:boruvka-round}
\end{figure}

\begin{complexitybox}
Each round scans all $m$ edges to find every component's cheapest outgoing edge and does $\Oh(m)$ union--find work: $\Oh(m\,\alpha(n))$ per round, or $\Oh(m)$ with a relabel-to-representatives bookkeeping. The component count at least halves each round, so there are $\le\lceil\log_2 n\rceil$ rounds. Total: $\Oh(m\log n)$. Boruvka's advantage is contraction: after each round you can rebuild the graph on the surviving supernodes (removing self-loops and keeping the lightest parallel edge), and \emph{that} shrinking is what the linear-time algorithms exploit.
\end{complexitybox}

\begin{aestheticbox}
All three algorithms apply the same theorem. Kruskal picks cuts globally cheapest-first; Prim fixes one side of the cut and grows it; Boruvka uses every component's cut at once. They are three readings of the cut property, and each is optimal in a different regime: Kruskal when edges are presorted or radix-sortable, Prim with a Fibonacci heap on dense graphs ($\Oh(m+n\log n)$), Boruvka when you need parallelism or contraction. The randomized linear-time algorithm in Section~\ref{ch06:linear} is, structurally, Boruvka steps glued to a sampling trick.
\end{aestheticbox}

\section{Reverse-delete: the cycle property as an algorithm}

Kruskal, Prim, and Boruvka all build \emph{up} from the empty forest using the cut property. The cycle property gives a fourth algorithm that tears \emph{down} from the full graph: sort the edges heaviest-first and delete each one whose removal does not disconnect the graph.

\begin{intuitionbox}
Start with every edge present. Walk the edges from most expensive to cheapest. Throw an edge away whenever you can do so without breaking connectivity -- if its endpoints are still joined by some other route, the edge was the heaviest rung of a cycle and the cycle property says no MST needs it. Stop when you have removed everything removable; what survives is the MST.
\end{intuitionbox}

\begin{algorithm}[H]
\caption{Reverse-Delete-MST}
\label{ch06:reverse-delete}
\KwIn{connected weighted graph $\Gph=(V,E,w)$}
\KwOut{MST edge set $T$}
$T\gets E$\;
sort $E$ so that $w(e_1)\ge w(e_2)\ge\dots\ge w(e_m)$\;
\ForEach{edge $e$ in nonincreasing weight order}{
  \If{$(V,T\setminus\set e)$ is still connected}{
    $T\gets T\setminus\set e$\;
  }
}
\Return $T$\;
\end{algorithm}

\begin{whybox}
\textbf{Why reverse-delete is correct.} When we consider $e$, it is deleted exactly when $T\setminus\set e$ is still connected, i.e.\ when $e$'s endpoints are joined by a path of remaining edges. That path plus $e$ is a cycle, and because we go heaviest-first, every edge still present is no heavier than $e$ -- so $e$ is a maximum-weight edge on that cycle. The cycle property (Theorem~\ref{ch06:cycle-property}) says such an edge is in no MST, so deleting it is safe. Conversely an edge we \emph{keep} is a bridge of the surviving graph at the moment we test it; a bridge lies on no cycle, so removing it would disconnect the graph and it must stay. At the end $T$ is connected, acyclic (any remaining cycle would have a heaviest edge we would have deleted), and spans -- an MST.
\end{whybox}

\begin{labbox}
Trace reverse-delete on $\Gph_0$. Edges heaviest-first:
$EF(9)$, $CF(8)$, $CE(7)$, $BC(6)$, $BE(5)$, $AB(4)$, $DE(3)$, $BD(2)$, $AD(1)$.

\begin{center}
\begin{tabular}{cll}
\toprule
edge & weight & action \\
\midrule
$EF$ & $9$ & delete ($F$ still reached via $CF$) \\
$CF$ & $8$ & \emph{keep} (deleting isolates $F$ -- bridge) \\
$CE$ & $7$ & delete ($C,E$ still joined via $B,D$) \\
$BC$ & $6$ & \emph{keep} (only link left to $C$) \\
$BE$ & $5$ & delete ($B,E$ still joined via $D$) \\
$AB$ & $4$ & delete ($A,B$ still joined via $D$) \\
$DE$ & $3$ & \emph{keep} (only link left to $E$) \\
$BD$ & $2$ & \emph{keep} (bridge) \\
$AD$ & $1$ & \emph{keep} (bridge) \\
\bottomrule
\end{tabular}
\end{center}

Survivors $\set{CF,BC,DE,BD,AD}$, weight $20$ -- the same MST again, this time obtained by deleting edges. The kept edges are precisely the bridges; the deleted edges are exactly the cycle-maxima Kruskal rejected. Naively the connectivity test is a DFS per edge, giving $\Oh(m(m+n))$; with a dynamic-connectivity structure it drops to $\Oh(m\log^2 n)$, which is why reverse-delete is valued mainly for its proof and is rarely run as written.
\end{labbox}

\section{Uniqueness of the MST}
\label{ch06:uniqueness}

\begin{theorem}[Distinct weights $\Rightarrow$ unique MST]
\label{ch06:unique}
If all edge weights of a connected graph $\Gph$ are pairwise distinct, then $\Gph$ has exactly one minimum spanning tree.
\end{theorem}

\begin{proof}
Suppose $T_1\ne T_2$ are both MSTs. Let $e$ be the minimum-weight edge in their symmetric difference $T_1\triangle T_2$ (it is nonzero, and weights are distinct, so $e$ is unique). Without loss of generality $e\in T_1\setminus T_2$. Adding $e$ to $T_2$ creates a unique cycle $\gamma$. Not all edges of $\gamma$ can be in $T_1$ (it is a tree, no cycle), so $\gamma$ contains some edge $f\in T_2\setminus T_1\subseteq T_1\triangle T_2$. By minimality of $e$ in the symmetric difference, $w(e)<w(f)$. Then $T_2-f+e$ is a spanning tree of weight $w(T_2)-w(f)+w(e)<w(T_2)$, contradicting that $T_2$ is an MST.
\end{proof}

The converse fails: a graph with ties can still have a unique MST (if the tied edges happen never to compete). The precise criterion is local.

\begin{proposition}[When ties create alternatives]
\label{ch06:tie-criterion}
$\Gph$ has more than one MST iff there is an MST $T$ and a non-tree edge $f\notin T$ such that $f$ has the \emph{same} weight as some edge $g$ on the tree path between $f$'s endpoints in $T$. Swapping $g$ for $f$ produces a different MST of equal weight.
\end{proposition}

\begin{whybox}
The forward direction is the exchange argument: any two MSTs differ by a sequence of equal-weight swaps. If $T\ne T'$ are MSTs, pick $f\in T'\setminus T$; the $T$-path between its ends contains an edge $g\notin T'$ on the cycle $f$ closes, and equal total weight forces $w(g)=w(f)$ (if $w(g)<w(f)$ then $T'$ could be improved; if $w(g)>w(f)$ then $T$ could be). The reverse direction is immediate: the swap yields a distinct equal-weight tree.
\end{whybox}

\begin{remark}
A clean way to detect uniqueness algorithmically: build any MST, then for each non-tree edge $f$ check whether the heaviest edge on the tree-path between its endpoints equals $w(f)$. If some such pair ties, the MST is not unique. The path-max queries are exactly the verification primitive of Section~\ref{ch06:verify}, so uniqueness testing is no harder than MST verification: $\Oh(m+n)$.
\end{remark}

\begin{labbox}
A worked tie. Take a triangle $A,B,C$ with $w(AB)=w(BC)=1$ and $w(AC)=2$, plus a pendant $w(CD)=3$. Two MSTs exist: $\set{AB,BC,CD}$ and $\set{AB,AC,CD}$? No -- the second uses $AC(2)$, heavier than $BC(1)$, so it is not minimum. The genuine tie needs equal competitors: set $w(AB)=w(BC)=w(AC)=1$. Now any two of the three triangle edges form a spanning tree of the triangle, all of weight $2$, so with the pendant $CD$ there are three MSTs: $\set{AB,BC,CD}$, $\set{AB,AC,CD}$, $\set{BC,AC,CD}$. Proposition~\ref{ch06:tie-criterion} flags this immediately: in $T=\set{AB,BC,CD}$ the non-tree edge $AC$ has weight $1$, equal to the heaviest edge ($AB$ or $BC$, both $1$) on the tree path $A\!-\!B\!-\!C$, so a swap exists. Distinct weights would have broken exactly this symmetry.
\end{labbox}

\begin{connectionbox}
The MST lives inside a \emph{matroid}: the spanning forests of a graph are the independent sets of the \emph{graphic matroid}, and the greedy algorithm is optimal for \emph{every} matroid (the matroid greedy theorem). The cut and cycle properties are the graphic specialization of a matroid's cocircuit and circuit elimination axioms. This is why MST greed works and, say, greedy on shortest paths or on the travelling salesman does not: those feasibility structures are not matroids. The companion linear-algebra book's treatment of the cycle space and cut space of a graph -- the kernel and row space of the incidence matrix -- is the same duality you met here as cycle property versus cut property.
\end{connectionbox}

\section{Verifying a spanning tree in linear time}
\label{ch06:verify}

A question that sounds easier than building an MST leads to the hardest result in the chapter: \emph{given} a spanning tree $T$, is it minimum? Recomputing the MST and comparing is too slow, because the randomized algorithm later in the chapter needs verification to run in \emph{linear} time.

\begin{focusbox}
\textbf{MST verification, the test.} A spanning tree $T$ is an MST iff for every non-tree edge $f=(u,v)$, $w(f)\ge\max\set{w(g):g\text{ on the }T\text{-path from }u\text{ to }v}$. Equivalently: no non-tree edge is lighter than the heaviest tree edge it would replace.
\end{focusbox}

\begin{whybox}
This is precisely the cycle property read both ways. If some $f$ is lighter than the heaviest edge $g$ on its tree path, the cycle $f+\text{path}$ has its maximum at $g$ (a tree edge), so by Theorem~\ref{ch06:cycle-property} $g$ is in no MST, hence $T$ is not minimum -- and indeed $T-g+f$ is lighter. Conversely, if the test passes for every $f$, then $T$ satisfies the cycle-optimality condition for every fundamental cycle, which is known to be equivalent to minimality (each non-tree edge is a maximum on its own fundamental cycle).
\end{whybox}

So verification reduces to: for each of the $m$ non-tree edges, compute the maximum-weight edge on a tree path -- the \emph{tree-path-maximum} problem. Done naively (walk each path) this is $\Oh(nm)$. In fact all $m$ queries together take $\Oh(m+n)$ time: Koml\'os (1985) proved the comparison bound, and Dixon--Rauch--Tarjan (1992) gave a linear-time implementation.

\begin{intuitionbox}
Koml\'os's insight: you do not answer the $m$ queries one at a time. You push all queries down into the tree at once and answer them in a single sweep, reusing comparisons across queries that share path-prefixes. The number of \emph{comparisons} needed to answer $m$ path-max queries on an $n$-node tree is $\Oh(n+m)$ -- provably, by an amortization that charges each comparison to a query or a tree edge. Turning that comparison bound into a linear-\emph{time} algorithm needs a Boruvka-style reduction to a balanced tree plus careful bit-packing of the answers; we sketch the structure rather than the full pointer-arithmetic.
\end{intuitionbox}

\begin{algorithm}[H]
\caption{MST-Verify (Koml\'os, structure)}
\label{ch06:verify-alg}
\KwIn{graph $\Gph$, spanning tree $T$}
\KwOut{\textsc{true} iff $T$ is an MST}
contract $T$ via Boruvka steps into a balanced tree $B$ of depth $\Oh(\log n)$ where every internal node has $\ge 2$ children\;
\ForEach{non-tree edge $f=(u,v)$}{
  route the query to the node $\mathrm{lca}_B(u,v)$\;
}
process all queries in one bottom-up sweep of $B$, maintaining for each query the running path-maximum, packing the per-edge comparison results into machine words\;
\lForEach{non-tree edge $f$}{\lIf{$w(f)<\mathrm{pathmax}(f)$}{\Return \textsc{false}}}
\Return \textsc{true}\;
\end{algorithm}

\begin{whybox}
Why a \emph{balanced} contracted tree? On a path-shaped tree of depth $n$, a single query could touch $n$ edges, and $m$ queries naively cost $nm$. Boruvka contraction replaces $T$ by a tree of height $\Oh(\log n)$ where each node has at least two children and the path-maxima of the original tree are recoverable from path-maxima in the contracted one. On the balanced tree each query path has length $\Oh(\log n)$, and Koml\'os's amortization brings the \emph{total} comparison count to $\Oh(n+m)$: each tree edge sees $\Oh(1)$ amortized comparisons because the branching factor $\ge 2$ geometrically bounds how many queries can ``stack up'' on any edge. Pack the comparison outcomes along each leaf-to-root path into $\Oh(\log n)$-bit words so that a query reads its answer in $\Oh(1)$ word operations.
\end{whybox}

\begin{complexitybox}
MST verification: $\Oh(m+n)$ time on a unit-cost RAM. Number of comparisons (Koml\'os bound): $\Oh(m+n)$, indeed at most $\approx m\log_2(2 + n/m)$ in the refined count. The randomized MST algorithm uses this linear verifier to cheaply discard the edges that provably cannot be in the MST.
\end{complexitybox}

\section{Randomized linear-time MST}
\label{ch06:linear}

Whether MST can be solved in deterministic linear time on a comparison RAM remains open. Karger, Klein and Tarjan (1995) gave a randomized algorithm running in $\Oh(m)$ expected time. Chazelle (2000) gave a deterministic $\Oh(m\,\alpha(m,n))$ algorithm, and Pettie and Ramachandran (2002) found an algorithm provably optimal up to constants. The randomized algorithm is the simplest of these to present in full.

The algorithm rests on one lemma. Sample a random subgraph, compute its minimum spanning \emph{forest}, and use that forest to certify that most edges are useless.

\begin{definition}[$F$-heavy and $F$-light]
Let $F$ be a forest in $\Gph$. An edge $f=(u,v)\in E$ is \emph{$F$-heavy} if $u,v$ lie in the same tree of $F$ and $w(f)$ exceeds the maximum weight on the $F$-path between them; otherwise $f$ is \emph{$F$-light}. By the cycle property, an $F$-heavy edge is in no MST of $\Gph$.
\end{definition}

Detecting which edges are $F$-heavy is exactly the verification primitive: $\Oh(m+n)$ for all edges at once.

\begin{lemma}[Sampling / DJP--KKT lemma]
\label{ch06:sampling}
Let $H$ be the subgraph obtained by including each edge of $\Gph$ independently with probability $p$, and let $F$ be the minimum spanning forest of $H$. Then the expected number of $F$-light edges of $\Gph$ is at most $n/p$, where $n=\abs V$.
\end{lemma}

\begin{whybox}
The proof considers edges in increasing weight order and runs Kruskal on $\Gph$ and on $H$ simultaneously. An edge is $F$-light exactly when, at the moment Kruskal would consider it, it would be \emph{accepted} into the forest of $H$ -- it connects two different $F$-components. Now flip the experiment: reveal each edge's coin only when Kruskal reaches it. An $F$-light edge that is accepted into $F$ is one of at most $n-1$ such acceptances; an $F$-light edge that is \emph{not} sampled still ``would have been accepted,'' and each such edge is sampled with probability $p$ independently, so the count of $F$-light edges is dominated by a sequence of coin flips with at most $n-1$ successes, whose expected length is $\le (n-1)/p < n/p$. So light edges are rare, because the random forest already connects most of the graph.
\end{whybox}

\begin{algorithm}[H]
\caption{KKT randomized MST (expected $\Oh(m)$)}
\label{ch06:kkt}
\KwIn{connected weighted graph $\Gph=(V,E,w)$}
\KwOut{MST edge set}
\If{$\abs E$ is small (base case)}{\Return MST by Boruvka/Prim directly}
run \textbf{3 Boruvka steps} on $\Gph$, contracting; collect the chosen edges in $T_0$\;
\tcp{now $\le n/8$ vertices remain; call the contracted graph $\Gph'$}
let $H$ be a random subgraph of $\Gph'$, each edge kept independently w.p.\ $1/2$\;
$F\gets$ MSF$(H)$ \tcp*{recursive call}
discard every $F$-heavy edge of $\Gph'$ \tcp*{linear verification}
$T_1\gets$ MSF of the surviving subgraph \tcp*{recursive call}
\Return $T_0\cup T_1$\;
\end{algorithm}

\begin{whybox}
\textbf{Why it is correct.} The three Boruvka steps add only safe edges (cut property), so $T_0$ is a subset of an MST, and contracting along them preserves the MST of the rest. Discarding $F$-heavy edges removes only edges that, by the cycle property, are in no MST -- so the surviving graph still contains the MST of $\Gph'$. The two recursive calls compute spanning forests, and the union $T_0\cup T_1$ is therefore exactly an MST of $\Gph$.

\textbf{Why it is linear in expectation.} Let $T(m,n)$ bound the expected time on $m$ edges, $n$ vertices. The Boruvka steps cost $\Oh(m)$ and leave $\le n/8$ vertices. The first recursion runs on $H$, which has $\le m/2$ edges in expectation on $\le n/8$ vertices. By the sampling lemma with $p=1/2$, the expected number of $F$-light survivors is $\le 2\cdot(n/8)=n/4$, so the second recursion runs on $\le n/4$ edges (only $F$-light edges survive). Verification and discarding are $\Oh(m)$. Writing $E[T(m,n)]\le c\,m + E[T(m/2,\,n/8)] + E[T(n/4,\,n/8)]$ and solving (the two subproblem edge-counts sum to $\le m/2 + n/4 \le 3m/4$ since $n\le m$ in a connected graph) gives $E[T(m,n)]=\Oh(m)$ by a standard ``shrinking total size'' induction.
\end{whybox}

\begin{complexitybox}
KKT: $\Oh(m)$ \emph{expected} time, and in fact $\Oh(m)$ with probability $1-e^{-\Om(m)}$. The three ingredients: Boruvka steps to shrink the vertex set geometrically, random sampling to estimate the MST (sampling lemma), and linear-time verification to prune $F$-heavy edges using that sample. Each ingredient reuses a result from earlier in this chapter.
\end{complexitybox}

\begin{aestheticbox}
The pieces fit together as follows: Boruvka steps contract vertices, random sampling and verification prune edges, and recursion solves the smaller subproblems. The cut property licenses the contraction, the cycle property licenses the deletion. Every main result of this chapter is used in the final algorithm, which is a good reason to build the theory on the cut and cycle properties.
\end{aestheticbox}

\section{Applications}

\begin{applicationbox}
\textbf{Network design.} The original motivation: lay the cheapest cabling, piping, or road network that connects every site. MST gives the cost-minimal connected backbone; redundancy is then added on top deliberately.
\end{applicationbox}

\begin{applicationbox}
\textbf{Clustering (single-linkage).} Build the MST, then delete its $k-1$ heaviest edges; the resulting $k$ components are exactly the single-linkage clusters. Because the MST orders merges by increasing dissimilarity, this is hierarchical agglomerative clustering done in $\Oh(m\log n)$ instead of the naive $\Oh(n^3)$. The deleted edges are the largest ``gaps'' between clusters.
\end{applicationbox}

\begin{applicationbox}
\textbf{Approximation algorithms.} The MST is a $2$-approximation skeleton for metric TSP: double the MST, shortcut the resulting Eulerian walk, and you get a tour of length $<2\cdot\mathrm{OPT}$. Christofides refines this to $3/2$ by adding a minimum matching on odd-degree vertices. The MST also lower-bounds the Steiner tree.
\end{applicationbox}

\begin{applicationbox}
\textbf{Image segmentation and graph-based vision.} Felzenszwalb--Huttenlocher segmentation is essentially Kruskal with a stopping rule, merging pixel regions in increasing edge-weight order. The MST encodes the ``most natural'' merges.
\end{applicationbox}

\begin{connectionbox}
\textbf{To the spectral and electrical material later in the book.} The number of spanning trees of $\Gph$ is $\det(\Lap_{ii})$, any cofactor of the graph Laplacian -- the Matrix-Tree Theorem. Random spanning trees (sampled by Wilson's algorithm or by electrical-network methods) connect MSTs to random walks and to the effective resistances you will meet in the Laplacian-solver chapters. The cut/cycle duality of this chapter is the discrete counterpart of decomposing edge space orthogonally into the cut space and cycle space of the incidence matrix.

While a minimum spanning tree minimizes the total weight of all edges across the entire network, it does not minimize the path distance between any individual pair of vertices. If you want the shortest route from a specific source to every other node, a global minimum tree will often lead you down circuitous detours. In Chapter~7, we shift from global network economy to point-to-point distance: we lay the foundations of Shortest Paths, explore edge relaxation, and develop Dijkstra's algorithm for nonnegative weights.
\end{connectionbox}

\section*{Exercises}
\addcontentsline{toc}{section}{Exercises}

\begin{exercise}[Cut property, by hand]
On the graph $\Gph_0$ of Figure~\ref{ch06:running-graph}, take the cut $S=\set{A,D,E}$. List every crossing edge and identify the minimum. Verify that this edge is in the MST of Figure~\ref{ch06:mst}, as the cut property predicts.
\end{exercise}

\begin{exercise}[Reverse-delete]
Write pseudocode for the reverse-delete algorithm: sort edges in \emph{decreasing} weight, and delete an edge whenever its removal keeps the graph connected. Prove it correct using the cycle property, and state its running time if connectivity is tested by a single DFS per candidate edge.
\end{exercise}

\begin{exercise}[Prim vs.\ Dijkstra]
Construct a small weighted graph on which Prim's tree and Dijkstra's shortest-path tree (from the same root) differ, and exhibit a vertex whose two parent edges differ. Explain the discrepancy in one sentence using the priorities each algorithm minimizes.
\end{exercise}

\begin{exercise}[Boruvka ties]
Give a $4$-vertex graph with repeated weights on which naive Boruvka (no tie-break) can add a cycle. Then show the lexicographic $(w,\text{id})$ tie-break prevents it. How many Boruvka rounds does $K_n$ with all-equal weights take under that tie-break?
\end{exercise}

\begin{exercise}[Second-best spanning tree]
Design an $\Oh(m\log n)$ algorithm to find the second-minimum spanning tree (the cheapest tree that is not the MST). Hint: it differs from the MST by exactly one swap; for each non-tree edge $f$, consider replacing the maximum tree edge on $f$'s fundamental cycle, and use the path-max primitive.
\end{exercise}

\begin{exercise}[Uniqueness test]
Prove that an MST is unique iff for every non-tree edge $f$, $w(f)$ is strictly greater than the maximum tree-path edge weight between $f$'s endpoints. Turn this into an $\Oh(m+n)$ test assuming an $\Oh(m+n)$ path-max oracle.
\end{exercise}

\begin{exercise}[Bottleneck spanning tree]
A \emph{minimum bottleneck spanning tree} minimizes its single \emph{heaviest} edge. Prove every MST is also a minimum bottleneck spanning tree, but the converse can fail. Give a counterexample for the converse.
\end{exercise}

\begin{exercise}[Edge weight changes]
Suppose one edge $e$'s weight changes. (a) If $e\in T$ and its weight \emph{decreases}, argue $T$ stays an MST. (b) If $e\notin T$ and its weight decreases below the path-max on its fundamental cycle, describe the single swap that repairs the MST. Give an $\Oh(n)$ update procedure for both cases.
\end{exercise}

\begin{exercise}[Sampling lemma sharpness]
In Lemma~\ref{ch06:sampling}, show the bound $n/p$ is tight up to constants by exhibiting a graph and a $p$ where the expected number of $F$-light edges is $\Om(n/p)$. Conclude why $p=1/2$ is a sensible choice in the KKT recursion.
\end{exercise}

\begin{exercise}[Implement KKT]
Implement Karger--Klein--Tarjan using your union--find and a path-max verifier. Empirically measure the number of edges entering each recursion level on random graphs and check it shrinks geometrically, matching the $\Oh(m)$ expected-time analysis. Where does the constant hide?
\end{exercise}

\begin{exercise}[Matroid greedy]
State the graphic matroid (ground set $E$, independent sets = acyclic edge subsets) and prove the exchange axiom. Explain in two sentences why the greedy algorithm optimal for matroids specializes to Kruskal.
\end{exercise}

\begin{exercise}[Boruvka contraction cost]
Show that after a Boruvka round you can rebuild the contracted multigraph -- relabel endpoints to component representatives, drop self-loops, keep one cheapest copy of each parallel edge -- in $\Oh(m)$ time. Conclude the standalone Boruvka algorithm runs in $\Oh(m\log n)$.
\end{exercise}
